\documentclass[11pt,a4paper]{article}

% Langue et polices (compiler avec XeLaTeX)
\usepackage{fontspec}
\usepackage[french]{babel}

\usepackage[margin=2.5cm]{geometry}
\usepackage{amsmath, amssymb}
\usepackage{graphicx}
\usepackage{enumitem}
\usepackage[hidelinks]{hyperref}

% Informations du devoir
\newcommand{\cours}{Nom du cours}
\newcommand{\devoir}{Devoir 1}
\newcommand{\etudiant}{Prénom Nom}
\newcommand{\professeur}{Nom du professeur}

% Question numérotée et environnement de réponse
\newcounter{question}
\newcommand{\question}[1]{%
  \stepcounter{question}%
  \section*{Question \thequestion\ifx\relax#1\relax\else: #1\fi}}
\newenvironment{reponse}{\par\noindent\textbf{Réponse.}\ }{\hfill$\blacksquare$\par}

\begin{document}

\noindent
\begin{minipage}[t]{0.6\textwidth}
  \textbf{\cours}\\
  \devoir
\end{minipage}%
\begin{minipage}[t]{0.4\textwidth}
  \raggedleft
  \etudiant\\
  Professeur : \professeur\\
  \today
\end{minipage}
\par\medskip\hrule\medskip

\question{Dérivées}
Calculer $\dfrac{d}{dx}\left(x^2 \sin x\right)$.

\begin{reponse}
  D’après la règle du produit,
  \[
    \frac{d}{dx}\left(x^2 \sin x\right) = 2x \sin x + x^2 \cos x.
  \]
\end{reponse}

\question{}
Énoncer deux propriétés de l’intégrale définie.

\begin{reponse}
  \begin{enumerate}[label=(\alph*)]
    \item Linéarité: $\int_a^b (\alpha f + \beta g) = \alpha \int_a^b f + \beta \int_a^b g$.
    \item Additivité: $\int_a^b f = \int_a^c f + \int_c^b f$.
  \end{enumerate}
\end{reponse}

\end{document}
